\documentclass[11pt]{letter}
\usepackage[utf8]{inputenc}
\usepackage[T1]{fontenc}
\usepackage[a4paper,top=2.4cm,bottom=2.4cm,left=2.5cm,right=2.5cm]{geometry}
\usepackage{lmodern}
\usepackage{microtype}
\usepackage{hyperref}
\hypersetup{colorlinks=true, urlcolor=blue, linkcolor=black}
\setlength{\parindent}{0pt}
\setlength{\parskip}{8pt}
\pagestyle{empty}

\begin{document}

\begin{flushright}
Md Rakib Hasan\\
Fermium Systems, Dhaka 1207, Bangladesh\\
Department of Soil, Water and Environment, University of Dhaka\\
\href{mailto:rakibhhridoy@fermium.systems}{rakibhhridoy@fermium.systems}
\end{flushright}

To the Editors,\\
\textit{Nature Communications}

\textbf{Re: Submission of an original research manuscript}

Dear Editors,

We are pleased to submit our manuscript, \textbf{``Global maps overstate the reliability
of blue-carbon credits,''} for consideration as an article in \textit{Nature
Communications}.

Mangrove soil and biomass carbon underpins a fast-growing market: global wall-to-wall
maps are now used to set baselines for blue-carbon credits and to target restoration. Yet
these maps are validated almost exclusively with random or spatially-pooled
cross-validation, which is inflated by spatial autocorrelation and silent on the question
a carbon buyer or a newly-surveyed coastline actually faces: can a model trained on some
regions predict an \emph{unsampled} one? We show, across a reproducible multi-delta
benchmark and 29 data-driven mangrove regions on six continents, that it cannot, and we
quantify what that means for the integrity of carbon-crediting baselines.

We believe the study is well suited to the broad readership of \textit{Nature
Communications} for three reasons:

\begin{itemize}
  \item \textbf{A general and consequential result.} Under leave-one-delta-out evaluation
  a model scoring $R^2=0.65$ in random cross-validation captures essentially no
  within-region variation out of region (median correlation $r\approx0$), and every
  held-out delta lies entirely outside the model's area of applicability, a finding robust
  to the applicability threshold and reproduced across 29 regions spanning six continents.
  A widely-used published soil-carbon product shows the same per-region pattern against
  in-situ cores, so the caution applies to maps already in use. For the 8\% of global
  mangrove area in deltas that lack in-situ cores, the stock that sets a crediting baseline
  cannot be pinned by a global model to better than a sixfold range, an uncertainty no
  current product reports.
  \item \textbf{A mechanism, a specificity test, and a constructive remedy.} The failure
  reflects concept shift, the environment--carbon relationship itself differing between
  deltas, rather than missing predictors, which adding geomorphic and tidal forcing does
  not remedy; and it is specific to blue carbon, as terrestrial soil carbon transfers far
  better between continents under the identical protocol. We then show that global and
  local data are complementary rather than substitutable: neither a global model alone nor
  a handful of local cores alone predicts within-delta carbon, but together, on the order
  of ten to twenty-five local cores, they recover most of the lost skill.
  \item \textbf{A reusable, uncertainty-aware framework.} We release a frozen multi-delta
  benchmark with leave-one-delta-out and leave-one-region-out splits, a three-tier
  validation protocol, and area-of-applicability and conformal-prediction layers that tell
  a practitioner where a global model can be trusted, portable to any coastal system by a
  bounding box.
\end{itemize}

We have been deliberately conservative in interpretation: we report median alongside mean
skill, separate level error from pattern error rather than leaning on a single composite
metric, and state the limitations of sparse in-situ labels and the soil-carbon-led scope
explicitly. The complete pipeline, pinned environment, and all derived data are openly
archived (Zenodo, \href{https://doi.org/10.5281/zenodo.20787021}{10.5281/zenodo.20787021})
to ensure full reproducibility.

\textbf{Declarations.} This manuscript is original, has not been published previously, and
is not under consideration elsewhere. The study involved no human participants or animals.
All primary datasets are public, and our code and derived data are archived to ensure
reproducibility. The author declares no competing interests.

Suitable reviewers may be drawn from the fields of blue-carbon biogeochemistry, spatial
machine learning for ecological mapping, and mangrove remote sensing; we are happy to
provide names on request.

Thank you for considering our manuscript. We look forward to your response.

Sincerely,

Md Rakib Hasan

\end{document}
